% -*- coding: utf-8 -*-
% !TEX program = xelatex
\documentclass[plain,noanswer,medmath]{../../template/jnuexam}
\usepackage{../../template/kysx}

\begin{document}


\examtitle{2017}{2}


\loadproblems{1}{2017P1}%载入试卷一


\makepart{选择题}{1～8小题，每小题4分，共32分}


\useproblem{1}{1}{1}%【同试卷一第一(1)题】


\begin{problem}
设二阶可导函数$f(x)$满足$f(1) = f(- 1) = 1$，$f(0) = - 1$，且$f''(x) > 0$，则\pickout{}
\begin{abcd}
\item $\int_{-1}^1 f(x)\dx > 0$．
\item $\int_{-1}^1 f(x)\dx < 0$．
\item $\int_{-1}^0 f(x)\dx > \int_0^1 f(x)\dx$．
\item $\int_{-1}^0 f(x)\dx < \int_0^1 f(x)\dx$．
\end{abcd}
\end{problem}


\begin{problem}
设数列$\left\{x_n \right\}$收敛，则\pickout{}
\begin{abcd}
\item 当$\lim_{n \to \infty} \sin x_n = 0$时，$\lim_{n \to \infty} x_n = 0$．
\item 当$\lim_{n \to \infty} (x_n + \sqrt{\left| x_n \right|}) = 0$时，$\lim_{n \to \infty} x_n = 0$．
\item 当$\lim_{n \to \infty} (x_n + x_n^2) = 0$时，$\lim_{n \to \infty} x_n = 0$．
\item 当$\lim_{n \to \infty} (x_n + \sin x_n) = 0$时，$\lim_{n \to \infty} x_n = 0$．
\end{abcd}
\end{problem}


\begin{problem}
微分方程$y'' - 4y' + 8y = \e^{2x}(1 + \cos 2x)$的特解可设为$y^* =$\pickout{}
\begin{abcd}
\item $A\e^{2x} + \e^{2x}\left(B\cos 2x + C\sin 2x \right)$．
\item $Ax \e^{2x} + \e^{2x}\left(B\cos 2x + C\sin 2x \right)$．
\item $A\e^{2x} + x\e^{2x}\left(B\cos 2x + C\sin 2x \right)$．
\item $Ax \e^{2x} + x\e^{2x}\left(B\cos 2x + C\sin 2x \right)$．
\end{abcd}
\end{problem}


\begin{problem}
设$f(x,y)$具有一阶偏导数，且对任意的$(x,y)$，
都有$\frac{\pd f(x,y)}{\pdx} > 0$，$\frac{\pd f(x,y)}{\pdy} < 0$，则\pickout{}
\begin{abcd}
\item $f(0,0) > f(1,1)$．
\item $f(0,0) < f(1,1)$．
\item $f(0,1) > f(1,0)$．
\item $f(0,1) < f(1,0)$．
\end{abcd}
\end{problem}


\useproblem{1}{1}{4}%【同试卷一第一(4)题】


\begin{problem}
设$A$为$3$阶矩阵，$P = (\alpha_1,\alpha_2,\alpha_3)$为可逆矩阵，使得$P^{-1} AP = \begin{pmatrix}
  0 & 0 & 0 \\
  0 & 1 & 0 \\
  0 & 0 & 2
\end{pmatrix}$，则$A(\alpha_1 + \alpha_2 + \alpha_3) =$\pickout{}
\begin{abcd}
\item $\alpha_1 + \alpha_2$．
\item $\alpha_2 + 2\alpha_3$．
\item $\alpha_2 + \alpha_3$．
\item $\alpha_1 + 2\alpha_2$．
\end{abcd}
\end{problem}


\useproblem{1}{1}{6}%【同试卷一第一(6)题】


\makepart{填空题}{9～14小题，每小题4分，共24分}


\begin{problem}
曲线$y = x\left(1 + \arcsin \frac{2}{x} \right)$的斜渐近线方程为 \fillin{}．
\end{problem}


\begin{problem}
设函数$y = y(x)$由参数方程$\begin{cases}
  x = t + \e^t \\
  y = \sin t
\end{cases}$确定，则$\left.\frac{\d^2y}{\dx^2}\right|_{t=0} =$ \fillin{}．
\end{problem}


\begin{problem}
$\int_0^{+\infty} \frac{\ln(1 + x)}{(1 + x)^2}\dx =$ \fillin{}．
\end{problem}


\begin{problem}
设函数$f(x,y)$具有一阶连续偏导数，且$\d f(x,y) = y\e^y \dx + x(1 + y) \e^y \dy$，$f(0,0) = 0$，
则$f(x,y)=$ \fillin{}．
\end{problem}


\begin{problem}
$\int_0^1 \dy \int_y^1 \frac{\tan x}{x}\dx =$ \fillin{}．
\end{problem}


\begin{problem}
设矩阵$A = \begin{pmatrix}
  4 & 1 & -2 \\
  1 & 2 & a \\
  3 & 1 & -1
\end{pmatrix}$的一个特征向量为$\begin{pmatrix}
  1 \\
  1 \\
  2
\end{pmatrix}$，则$a =$ \fillin{}．
\end{problem}


\makepart{解答题}{15～23小题，共94分}


\begin{problem}[points=10]
求极限$\lim_{x \to 0^+} \frac{\int_0^x \sqrt{x - t} \e^t \dt}{\sqrt{x^3}}$．
\end{problem}


\useproblem[points=10]{1}{3}{15}%【同试卷一第三(15)题】


\useproblem[points=10]{1}{3}{16}%【同试卷一第三(16)题】


\useproblem[points=10]{1}{3}{17}%【同试卷一第三(17)题】


\useproblem[points=11]{1}{3}{18}%【同试卷一第三(18)题】


\begin{problem}[points=11]
已知平面区域$D = \left\{(x,y)|x^2 + y^2 \le 2y \right\}$，计算二重积分$\iint_D (x + 1)^2 \dx\dy$．
\end{problem}


\begin{problem}[points=11]
设$y(x)$是区间$\left(0,\frac{3}{2} \right)$内的可导函数，且$y(1) = 0$．
点$P$是曲线$L:y = y(x)$上的任意一点．$L$在$P$处的切线与$y$轴相交于点$\left(0,Y_p\right)$，
法线与$x$轴相交于点$\left(X_p,0 \right)$，若$X_p = Y_p$，求$L$上点的坐标$(x,y)$满足的方程．
\end{problem}


\useproblem[points=11]{1}{3}{20}%【同试卷一第三(20)题】


\useproblem[points=11]{1}{3}{21}%【同试卷一第三(21)题】


\end{document}
